\medskip
\paragraph{Formal Lean statement.}
\begin{verbatim}
theorem prop_primary_box :
    ∃ c K C : ℝ, 0 < c ∧ 0 < K ∧ 0 < C ∧
      ∀ n : ℕ, 3 ≤ n →
      ∀ t : ℕ, (t : ℝ) ≥ C * ↑n ^ (3 : Nat) →
        |primaryCoreContribution n t - paperMainTerm n t|
          ≤ (K * (n : ℝ) ^ (2 : Nat) / (t : ℝ)
              + K * ((n : ℝ) ^ (5 / 2 : ℝ) / (t : ℝ) ^ (3 / 2 : ℝ))
              + K * (n : ℝ) ^ (6 : Nat) / (t : ℝ) ^ (2 : Nat)
              + K * Real.exp (-(c * (n : ℝ) ^ (2 : Nat))))
            * paperMainScale n t
\end{verbatim}
\begin{proof}[Proof sketch]
Assemble the short bridge chain from exact to corrected, corrected to quartic, quartic to cubic, and cubic to Gaussian scale. Then combine the second-order cubic approximation, the cubic correction tail, and the core-mass transport bound into one final estimate.
\end{proof}
\paragraph{Role in the route.}
It depends directly on primaryCoreContribution, exactCore\_correctedCore\_gap\_abs\_le\_gaussianF\_short, correctedCore\_quarticCore\_gap\_abs\_le\_gaussianF\_short, quarticCore\_cubicCore\_gap\_abs\_le\_gaussianF\_short, cubic\_core\_second\_order\_gap\_uniform, cubicT\_sq\_core\_gaussian\_correction\_tail\_le, coreMass\_gap\_le\_exp\_gaussianF, gaussianScale\_eq\_texPrefactor\_mul\_gaussianF. It is used directly by normalizedCount\_asymptotic, thm\_main\_intro.
\paragraph{Source location.}
\texttt{RequestProject/HadamardCn3Asymptotics.lean:57}
\medskip
